\documentclass[12pt,a4paper,openany,oneside]{book}

% --- Paquetes ---
\usepackage[utf8]{inputenc}
\usepackage[spanish, es-noshorthands]{babel} % Desactiva shorthands conflictivos (>, <, etc.)
\usepackage{amsmath, amsfonts, amssymb}
\usepackage{graphicx}
\usepackage{geometry}
\usepackage{hyperref}
\usepackage{booktabs}
\usepackage{fancyhdr}
\usepackage{listings}
\usepackage{xcolor}
\usepackage{tikz}
\usetikzlibrary{arrows.meta, positioning, shapes.geometric}

\geometry{margin=2.5cm}

% --- Colores y estilos para código Python ---
\definecolor{codegreen}{rgb}{0,0.6,0}
\definecolor{codegray}{rgb}{0.5,0.5,0.5}
\definecolor{codepurple}{rgb}{0.58,0,0.82}
\definecolor{backcolour}{rgb}{0.95,0.95,0.92}

\lstset{
    language=Python,
    backgroundcolor=\color{backcolour},   
    commentstyle=\color{codegreen},
    keywordstyle=\color{blue},
    numberstyle=\tiny\color{codegray},
    stringstyle=\color{codepurple},
    basicstyle=\ttfamily\small,
    breaklines=true,
    frame=single,
    showstringspaces=false
}

% --- Estilo de cabecera y pie ---
\pagestyle{fancy}
\fancyhf{}
\renewcommand{\headrulewidth}{0.4pt}
\renewcommand{\footrulewidth}{0.4pt}
\fancyhead[L]{\small Carlos César Sánchez Coronel}
\fancyhead[R]{\small Sesión 1: El Lenguaje de la IA}
\fancyfoot[L]{\small Carlos César Sánchez Coronel}
\fancyfoot[C]{\thepage}

% --- Portada ---
\title{
    \textbf{Sesión 1} \\ 
    \Huge \textbf{EL LENGUAJE DE LA IA} \\
    \Large Datos, Funciones y Python \\
    \normalsize Fundamentos matemáticos para inteligencia artificial
}
\author{
    \small Por \\ \vspace{0.2cm}
    \Large \textsc{Carlos César Sánchez Coronel}
}
\date{2026}

\begin{document}

\maketitle
\tableofcontents
\addcontentsline{toc}{chapter}{Índice general}

% ------------------------------------------------------------------
% CAPÍTULO 1: SESIÓN 1 - EL LENGUAJE DE LA IA
% ------------------------------------------------------------------
\chapter{El Lenguaje de la IA: Datos, Funciones y Python}

En esta primera sesión, sentaremos las bases conceptuales y prácticas que nos acompañarán durante todo el curso. Exploraremos cómo los datos se organizan matemáticamente, qué tipos de funciones son el ``ADN'' de los algoritmos de inteligencia artificial, y cómo Python (con sus librerías) se convierte en nuestro laboratorio para experimentar. El objetivo es que al finalizar, tengas un mapa mental claro de las herramientas matemáticas y su relación con los modelos que estudiarás en cursos posteriores.

\section{Conceptos Fundamentales}

\subsection{Conjuntos numéricos y su representación en IA}
Los datos en inteligencia artificial se organizan en estructuras matemáticas que permiten operaciones eficientes. Estas son las más importantes:

\begin{itemize}
    \item \textbf{Escalar}: un número real. Representa una magnitud individual, como una temperatura, un peso o un bias (sesgo) en una neurona.
    \begin{equation*}
        x = 5.2 \quad \text{(escalar)}
    \end{equation*}
    
    \item \textbf{Vector}: lista ordenada de números. Cada punto en un espacio de características se representa como un vector.
    \begin{equation*}
        \mathbf{x} = \begin{bmatrix} x_1 & x_2 & \dots & x_n \end{bmatrix} \quad \text{ej. } \mathbf{x} = [2.5, 1.8, 3.0]
    \end{equation*}
    
    \item \textbf{Matriz}: tabla bidimensional de números. Una imagen en escala de grises es una matriz donde cada elemento es un píxel.
    \begin{equation*}
        \mathbf{X} = \begin{bmatrix}
            x_{11} & x_{12} & \dots & x_{1m} \\
            x_{21} & x_{22} & \dots & x_{2m} \\
            \vdots & \vdots & \ddots & \vdots \\
            x_{n1} & x_{n2} & \dots & x_{nm}
        \end{bmatrix}
    \end{equation*}
    
    \item \textbf{Tensor}: generalización a más de dos dimensiones. Un lote de imágenes a color tiene dimensiones: [batch, alto, ancho, canales].
    \begin{equation*}
        \mathcal{T} \in \mathbb{R}^{B \times H \times W \times C}
    \end{equation*}
\end{itemize}

\begin{center}
\begin{tikzpicture}[scale=0.8]
% Escalar
\draw[fill=blue!20] (0,1) circle (0.5);
\node at (0,1) {5.2};
\node at (0,1.8) {Escalar};

% Vector
\draw[fill=green!20] (3,1) rectangle (5,2);
\node at (4,1.5) {[2.5, 1.8, 3.0]};
\node at (4,2.3) {Vector (1D)};

% Matriz
\draw[fill=red!20] (0,-2) rectangle (3,-0.5);
\foreach \i in {0.3,0.9,1.5,2.1} {
    \draw (\i,-1.2) -- (\i,-0.5);
}
\foreach \j in {-1.8,-1.2,-0.5} {
    \draw (0,\j) -- (3,\j);
}
\node at (1.5,-1) {$p_{11}$};
\node at (2.1,-1) {$p_{12}$};
\node at (0.9,-1.6) {$p_{21}$};
\node at (1.5,-1.6) {$p_{22}$};
\node at (3.5,-1.25) {Matriz (2D)};

% Tensor
\draw[fill=yellow!20] (5,-3) rectangle (9,0);
\draw (5,-1.5) -- (9,-1.5);
\draw (5,-2.5) -- (9,-2.5);
\draw (7,-3) -- (7,0);
\node at (7,-0.5) {B};
\node at (8,-1.8) {H};
\node at (6,-2.3) {W};
\node at (7.5,-3.3) {Tensor 4D};
\end{tikzpicture}
\captionof{figure}{Representación de escalares, vectores, matrices y tensores.}
\end{center}

\subsection{Funciones como ``cajas negras''}
Una función $f: X \rightarrow Y$ asigna a cada elemento de un conjunto de entrada (dominio) un elemento de un conjunto de salida (rango). En IA, los modelos son funciones complejas que aprendemos a partir de datos.

\begin{equation*}
f(x) = y \quad \text{o en forma vectorial} \quad \mathbf{y} = f(\mathbf{x})
\end{equation*}

Por ejemplo, una red neuronal que clasifica imágenes es una función que toma un tensor de píxeles y devuelve una probabilidad por clase.

\textbf{Tipos fundamentales de funciones:}

\begin{itemize}
    \item \textbf{Lineales}: $f(x) = mx + b$. Representan relaciones de proporcionalidad directa. En IA, son la base de la regresión lineal y de las capas fully-connected.
    \begin{center}
    \begin{tikzpicture}[scale=0.6]
    \draw[->] (-2,0) -- (2,0) node[right] {$x$};
    \draw[->] (0,-2) -- (0,2) node[above] {$f(x)$};
    \draw[blue, thick] (-1.5,-1.5) -- (1.5,1.5);
    \node at (1,1.5) {$f(x)=x$};
    \end{tikzpicture}
    \end{center}
    
    \item \textbf{Polinomiales}: $f(x) = a_n x^n + \dots + a_0$. Útiles para entender overfitting y para aproximar funciones complejas.
    \begin{equation*}
    f(x) = x^2 - 2x + 1
    \end{equation*}
    
    \item \textbf{Exponenciales}: $f(x) = e^x$, $f(x) = a^x$. Aparecen en softmax, en funciones de pérdida como cross-entropy, y en modelos de crecimiento.
    \begin{equation*}
    e^x = \sum_{n=0}^{\infty} \frac{x^n}{n!}
    \end{equation*}
    
    \item \textbf{Logarítmicas}: $f(x) = \ln x$. Son la inversa de las exponenciales y se usan en entropía y verosimilitud.
    
    \item \textbf{Trigonométricas}: $\sin x$, $\cos x$. Claves en procesamiento de señales y en positional encodings de transformers.
    \begin{equation*}
    \text{PE}(pos, 2i) = \sin\left(\frac{pos}{10000^{2i/d}}\right)
    \end{equation*}
\end{itemize}

\section{Funciones de Activación (Deep Learning)}
Son funciones no lineales que se aplican a la salida de una neurona. Permiten que la red aprenda relaciones complejas. A continuación, las más importantes con sus ecuaciones, propiedades y aplicaciones.

\subsection{Sigmoide}
\begin{equation*}
\sigma(x) = \frac{1}{1 + e^{-x}}
\end{equation*}
\begin{itemize}
    \item Dominio: $\mathbb{R}$
    \item Rango: $(0,1)$
    \item Propiedades: derivable, saturación para $|x|$ grande (gradientes pequeños).
    \item Uso: capa final en clasificación binaria (probabilidad).
\end{itemize}

\subsection{Tanh}
\begin{equation*}
\tanh(x) = \frac{e^x - e^{-x}}{e^x + e^{-x}}
\end{equation*}
\begin{itemize}
    \item Dominio: $\mathbb{R}$
    \item Rango: $(-1,1)$
    \item Propiedades: centrada en cero, saturación similar a sigmoide.
    \item Uso: en RNN clásicas y como alternativa a sigmoide.
\end{itemize}

\subsection{ReLU (Rectified Linear Unit)}
\begin{equation*}
\text{ReLU}(x) = \max(0, x)
\end{equation*}
\begin{itemize}
    \item Dominio: $\mathbb{R}$
    \item Rango: $[0, \infty)$
    \item Propiedades: no satura en la región positiva, computacionalmente eficiente. Puede morir (neuronas muertas si siempre negativas).
    \item Uso: la más común en capas ocultas de CNNs y MLPs.
\end{itemize}

\subsection{Leaky ReLU}
\begin{equation*}
\text{LReLU}(x) = \max(\alpha x, x) \quad \text{con } \alpha \approx 0.01
\end{equation*}
Soluciona el problema de neuronas muertas permitiendo un pequeño gradiente cuando $x<0$.

\subsection{ELU (Exponential Linear Unit)}
\begin{equation*}
\text{ELU}(x) = \begin{cases} x & x \ge 0 \\ \alpha(e^x - 1) & x < 0 \end{cases}
\end{equation*}
\begin{itemize}
    \item Propiedades: suave para $x<0$, media de activaciones cercana a cero.
\end{itemize}

\subsection{Softmax}
\begin{equation*}
\text{softmax}(\mathbf{z})_i = \frac{e^{z_i}}{\sum_{j=1}^{K} e^{z_j}}
\end{equation*}
\begin{itemize}
    \item Convierte un vector en una distribución de probabilidad (suma 1).
    \item Uso: capa final en clasificación multiclase.
\end{itemize}

\subsection{GELU (Gaussian Error Linear Unit)}
\begin{equation*}
\text{GELU}(x) = x \Phi(x) \quad \text{donde } \Phi \text{ es la CDF normal estándar}
\end{equation*}
\begin{itemize}
    \item Propiedades: combina ReLU y dropout, usada en transformers como BERT y GPT.
\end{itemize}

\subsection{Swish / SiLU}
\begin{equation*}
\text{Swish}(x) = x \cdot \sigma(x)
\end{equation*}
\begin{itemize}
    \item Propiedades: no monótona, suave, usada en modelos eficientes como EfficientNet.
\end{itemize}

\begin{center}
\begin{tikzpicture}[scale=0.8]
\begin{axis}[
    title={Funciones de activación},
    xlabel=$x$,
    ylabel=$f(x)$,
    domain=-4:4,
    samples=100,
    legend pos=north west,
    grid=major,
    width=12cm,
    height=8cm
]
\addplot[blue] {1/(1+exp(-x))};
\addlegendentry{Sigmoide}
\addplot[red] {tanh(x)};
\addlegendentry{Tanh}
\addplot[green] {max(0,x)};
\addlegendentry{ReLU}
\addplot[orange] {x*(1/(1+exp(-x)))};
\addlegendentry{Swish}
\end{axis}
\end{tikzpicture}
\captionof{figure}{Gráficas de funciones de activación comunes.}
\end{center}

\section{Funciones de Pérdida (Loss Functions)}
Miden cuán lejos está la predicción del modelo del valor real. Son la base del entrenamiento.

\subsection{Error Cuadrático Medio (MSE)}
\begin{equation*}
\text{MSE} = \frac{1}{n}\sum_{i=1}^{n} (y_i - \hat{y}_i)^2
\end{equation*}
\begin{itemize}
    \item Uso: regresión.
    \item Propiedades: diferenciable, sensible a outliers.
\end{itemize}

\subsection{Error Absoluto Medio (MAE)}
\begin{equation*}
\text{MAE} = \frac{1}{n}\sum_{i=1}^{n} |y_i - \hat{y}_i|
\end{equation*}
\begin{itemize}
    \item Uso: regresión robusta.
    \item Propiedades: menos sensible a outliers que MSE.
\end{itemize}

\subsection{Cross-Entropy (Entropía Cruzada)}
Para clasificación binaria:
\begin{equation*}
\text{BCE} = -\frac{1}{n}\sum_{i=1}^{n} \left[ y_i \log(\hat{y}_i) + (1-y_i)\log(1-\hat{y}_i) \right]
\end{equation*}
Para clasificación multiclase:
\begin{equation*}
\text{CE} = -\sum_{i=1}^{n} \sum_{c=1}^{K} y_{i,c} \log(\hat{y}_{i,c})
\end{equation*}
\begin{itemize}
    \item Uso: clasificación.
    \item Propiedades: el logaritmo castiga fuertemente predicciones erróneas con alta confianza.
\end{itemize}

\subsection{Hinge Loss}
\begin{equation*}
\text{Hinge} = \max(0, 1 - y_i \hat{y}_i) \quad \text{con } y_i \in \{-1,1\}
\end{equation*}
\begin{itemize}
    \item Uso: SVM (máquinas de soporte vectorial).
\end{itemize}

\subsection{Focal Loss}
\begin{equation*}
\text{FL} = -(1 - \hat{y}_i)^\gamma \log(\hat{y}_i)
\end{equation*}
\begin{itemize}
    \item Uso: problemas de clases desbalanceadas (detección de objetos).
\end{itemize}

\subsection{Contrastive Loss / Triplet Loss}
\begin{equation*}
\mathcal{L} = \max(0, d(a,p) - d(a,n) + \alpha)
\end{equation*}
\begin{itemize}
    \item Uso: aprendizaje de representaciones (face recognition, embeddings).
\end{itemize}

\section{Funciones de Similitud y Distancia}
Miden la cercanía entre vectores.

\subsection{Producto punto}
\begin{equation*}
\langle \mathbf{u}, \mathbf{v} \rangle = \sum_{i=1}^{n} u_i v_i
\end{equation*}
\begin{itemize}
    \item Uso: base de la atención en transformers, medida de similitud no normalizada.
\end{itemize}

\subsection{Similitud de Coseno}
\begin{equation*}
\text{cos}(\theta) = \frac{\mathbf{u} \cdot \mathbf{v}}{\|\mathbf{u}\| \|\mathbf{v}\|}
\end{equation*}
\begin{itemize}
    \item Uso: comparación de embeddings, motores de búsqueda semántica.
\end{itemize}

\subsection{Distancia Euclidiana (L2)}
\begin{equation*}
d(\mathbf{u},\mathbf{v}) = \sqrt{\sum_{i=1}^{n} (u_i - v_i)^2}
\end{equation*}
\begin{itemize}
    \item Uso: KNN, K-means, clustering.
\end{itemize}

\subsection{Distancia Manhattan (L1)}
\begin{equation*}
d(\mathbf{u},\mathbf{v}) = \sum_{i=1}^{n} |u_i - v_i|
\end{equation*}
\begin{itemize}
    \item Uso: robusta, usada en algunos algoritmos de recomendación.
\end{itemize}

\subsection{Kernel RBF}
\begin{equation*}
K(\mathbf{u},\mathbf{v}) = \exp\left(-\gamma \|\mathbf{u}-\mathbf{v}\|^2\right)
\end{equation*}
\begin{itemize}
    \item Uso: SVM, procesos Gaussianos.
\end{itemize}

\section{Lógica Proposicional y su Relación con la IA}
La lógica proposicional maneja valores de verdad (True/False) y conectores lógicos (AND, OR, NOT). En programación, estos se usan en condiciones. En IA, aparecen en:

\begin{itemize}
    \item \textbf{Árboles de decisión}: cada nodo interno es una condición lógica sobre una característica.
    \begin{center}
    \begin{tikzpicture}[
        every node/.style={draw, rectangle, rounded corners, align=center},
        level distance=1.5cm,
        sibling distance=3cm
    ]
    \node {¿Edad > 18?}
        child { node {¿Género = M?} 
            child { node {Pred: 0} }
            child { node {Pred: 1} }
        }
        child { node {¿Ingreso > 50k?}
            child { node {Pred: 1} }
            child { node {Pred: 0} }
        };
    \end{tikzpicture}
    \end{center}
    \item \textbf{Sistemas basados en reglas}: motores de inferencia usan reglas lógicas.
    \item \textbf{Programación lógica}: Prolog, por ejemplo.
    \item \textbf{Verificación de modelos}: en sistemas críticos.
\end{itemize}

\section{Notación de Sumatorias y Productorias}
Son herramientas para expresar sumas y productos de muchos términos de forma compacta.

\subsection{Sumatoria}
\begin{equation*}
\sum_{i=1}^{n} x_i = x_1 + x_2 + \dots + x_n
\end{equation*}
Ejemplo en una neurona: la suma ponderada de las entradas.
\begin{equation*}
z = \sum_{i=1}^{n} w_i x_i + b
\end{equation*}

\subsection{Productoria}
\begin{equation*}
\prod_{i=1}^{n} x_i = x_1 \cdot x_2 \cdots x_n
\end{equation*}
Ejemplo: la probabilidad conjunta de eventos independientes.
\begin{equation*}
P(X_1, \dots, X_n) = \prod_{i=1}^{n} P(X_i)
\end{equation*}

\section{Python: Librerías, Frameworks y Herramientas}
Python es el lenguaje estándar en IA por su ecosistema. Aquí las librerías esenciales.

\subsection{Librerías fundamentales}
\begin{itemize}
    \item \textbf{NumPy}: manipulación de tensores, operaciones vectorizadas, álgebra lineal.
    \item \textbf{SciPy}: funciones científicas (integración, optimización, señales).
    \item \textbf{Matplotlib} / \textbf{Seaborn}: visualización.
    \item \textbf{SymPy}: cálculo simbólico.
\end{itemize}

\subsection{Librerías de machine learning}
\begin{itemize}
    \item \textbf{Scikit-learn}: algoritmos clásicos (regresión, clasificación, clustering, PCA, KNN, SVM).
    \item \textbf{XGBoost}, \textbf{LightGBM}, \textbf{CatBoost}: gradient boosting.
\end{itemize}

\subsection{Frameworks de deep learning}
\begin{itemize}
    \item \textbf{TensorFlow / Keras}: producción, alta API.
    \item \textbf{PyTorch}: investigación, diferenciación automática dinámica.
    \item \textbf{JAX}: composición de funciones, aceleración.
\end{itemize}

\subsection{Librerías especializadas}
\begin{itemize}
    \item \textbf{NLP}: Hugging Face Transformers, NLTK, spaCy, Gensim.
    \item \textbf{Visión}: OpenCV, PIL, torchvision.
    \item \textbf{Audio}: librosa, torchaudio.
    \item \textbf{Grafos}: NetworkX, PyTorch Geometric.
    \item \textbf{Probabilidad}: PyMC3, TensorFlow Probability.
\end{itemize}

\subsection{Entornos y herramientas}
\begin{itemize}
    \item \textbf{Jupyter Notebook / JupyterLab}: experimentación interactiva.
    \item \textbf{Google Colab}: GPU/TPU gratuitas.
    \item \textbf{VS Code}, \textbf{PyCharm}: editores.
    \item \textbf{Docker}: reproducibilidad.
    \item \textbf{Git}: control de versiones.
\end{itemize}

\section{Ejemplo Práctico en Python}
Visualicemos algunas funciones de activación y pérdida con NumPy y Matplotlib.

\begin{lstlisting}[language=Python]
import numpy as np
import matplotlib.pyplot as plt

# Funciones de activación
def sigmoid(x):
    return 1 / (1 + np.exp(-x))

def tanh(x):
    return np.tanh(x)

def relu(x):
    return np.maximum(0, x)

def softmax(x):
    e_x = np.exp(x - np.max(x))
    return e_x / e_x.sum()

# Visualización
x = np.linspace(-5, 5, 100)
plt.figure(figsize=(12, 4))
plt.subplot(1, 3, 1)
plt.plot(x, sigmoid(x), label='Sigmoide')
plt.plot(x, tanh(x), label='Tanh')
plt.plot(x, relu(x), label='ReLU')
plt.legend()
plt.title('Funciones de activación')
plt.grid(True)

# Ejemplo de producto punto y similitud de coseno
u = np.array([1, 2, 3])
v = np.array([4, 5, 6])
dot = np.dot(u, v)
norm_u = np.linalg.norm(u)
norm_v = np.linalg.norm(v)
cos_sim = dot / (norm_u * norm_v)
print(f'Producto punto: {dot}')
print(f'Similitud de coseno: {cos_sim:.4f}')

# Softmax
z = np.array([2.0, 1.0, 0.1])
probs = softmax(z)
print(f'Softmax: {probs}, suma: {probs.sum():.2f}')
plt.subplot(1, 3, 2)
plt.bar(range(len(probs)), probs)
plt.title('Distribución softmax')
plt.xticks(range(len(probs)), ['Clase 0', 'Clase 1', 'Clase 2'])

# MSE
y_true = np.array([3, -0.5, 2, 7])
y_pred = np.array([2.5, 0.0, 2, 8])
mse = np.mean((y_true - y_pred)**2)
print(f'MSE: {mse:.4f}')
plt.subplot(1, 3, 3)
plt.scatter(y_true, y_pred)
plt.plot([-1, 8], [-1, 8], 'r--')
plt.xlabel('Valor real')
plt.ylabel('Predicción')
plt.title('MSE')
plt.grid(True)
plt.tight_layout()
plt.show()
\end{lstlisting}

\section{Diagrama de una Neurona Artificial}
Una neurona artificial calcula una combinación lineal de las entradas más un sesgo, y luego aplica una función no lineal $f$:
\begin{equation*}
y = f\left( \sum_{i=1}^{n} w_i x_i + b \right)
\end{equation*}

\begin{center}
\begin{tikzpicture}[
    neuron/.style={circle, draw, minimum size=0.8cm},
    weight/.style={font=\small},
    sum/.style={circle, draw, minimum size=0.6cm, inner sep=1pt},
    act/.style={rectangle, draw, rounded corners, minimum size=0.8cm},
    >=Stealth
]

% Entradas
\node[neuron] (x1) at (0,0) {$x_1$};
\node[neuron] (x2) at (0,-1.2) {$x_2$};
\node[neuron] (x3) at (0,-2.4) {$x_3$};

% Pesos (etiquetas)
\node[weight] at (1.2, -0.2) {$w_1$};
\node[weight] at (1.2, -1.4) {$w_2$};
\node[weight] at (1.2, -2.6) {$w_3$};

% Sumador
\node[sum] (sum) at (3, -1.2) {$\sum$};

% Bias
\node[weight] (bias) at (3, 0.2) {$+b$};

% Activación
\node[act] (act) at (5, -1.2) {$f$};

% Salida
\node[neuron] (y) at (7, -1.2) {$y$};

% Conexiones de entrada
\foreach \i in {1,2,3} {
    \draw[->] (x\i) -- (sum);
}

% Conexiones internas
\draw[->] (sum) -- (act);
\draw[->] (bias) -- (sum);
\draw[->] (act) -- (y);

\end{tikzpicture}
\captionof{figure}{Estructura de una neurona artificial con tres entradas.}
\end{center}

\section{Conclusión}
En esta sesión hemos construido un mapa conceptual que relaciona:
\begin{itemize}
    \item Las estructuras de datos fundamentales (escalares, vectores, matrices, tensores).
    \item Las funciones matemáticas que aparecen una y otra vez en IA (activación, pérdida, similitud).
    \item Los algoritmos y modelos donde se aplican (redes neuronales, SVM, clustering, transformers).
    \item Las herramientas de Python que usaremos para implementarlos.
\end{itemize}
Este conocimiento te servirá de brújula durante el resto del curso. ¡Bienvenido al fascinante mundo de las matemáticas para la inteligencia artificial!

\end{document}